\documentclass[12pt]{article}

\usepackage{fouriernc}
\usepackage{amsmath}

\usepackage{tikz}

\newcommand{\contradiction}{%
\begin{tikzpicture}[rotate=45,x=0.5ex,y=0.5ex]
\draw[line width=.2ex] (0,2) -- (3,2) (0,1) -- (3,1) (1,3) -- (1,0) (2,3) -- (2,0);
\end{tikzpicture}
}

\begin{document}

\section*{Problem-225}
\subsection*{Problem Statement}
If in the equation
\[
	ax^2 + bx + c = 0
\]
the coefficient $a$ is equal to zero then the equation takes the form
\[
	bx + c = 0
\]
and has the unique solution
\[
	x = -\frac{c}{b}
\]

Strictly speaking, the last sentence is wrong; when $b=0$ the quotient $\frac{c}{b}$ is undefined. How are we to correct this error?

\subsection*{Solution}
We should require that in the equation $ax^2+bx+c = 0$, both $a$ and $b$ cannot be zero.

\section*{Problem-227}
\subsection*{Problem Statement}
Prove that $\sqrt{3}$ is irrational.

\subsection*{Solution}
The statement `$\sqrt{3}$ is irrational' is equivalent to the statement: $\sqrt{3} \neq \frac{a}{b}$ for any integer $a$ and $b$. We shall make use of \textbf{proof by contradiction}, therefore, we shall assume that $\sqrt{3} = \frac{a}{b}$ and that should lead to a contradiction, completing the proof.

We shall consider four different cases for the pair $(a, b)$ and we shall derive a contradiction for each.
\begin{enumerate}
\item $a$ and $b$ both are odd integer. Say $a = 2m+1$ and $b = 2n+1$, where $m$ and $n$ are integers. According to our assumption
\begin{equation*}
	\begin{aligned}
		\sqrt{3} &= \frac{2m+1}{2n+1}\\
		3        &= \frac{ (2m+1)^2 }{ (2n+1)^2 }\\
		3\ (2n+1)^2 &= (2m+1)^2\\
		3\ (4n^2 + 4n + 1) &= 4m^2 + 4m + 1\\
		12n^2 + 12n + 3 &= 4m^2 + 4m + 1\\
		12n^2 + 12n + 2 &= 4m^2 + 4m\\
		6n^2 + 6n + 1 &= 2m^2 + 2m
	\end{aligned}
\end{equation*}
On the left side we have an odd integer and on the right side we have an even integer---a contradiction. \contradiction

\item $a$ is an odd integer and $b$ is an even integer. Say $a = 2m+1$ and $b = 2n$, where $m$ and $n$ are integers. According to our assumption
\begin{equation*}
	\begin{aligned}
		\sqrt{3} &= \frac{ 2m+1 }{ 2n }\\
		3 &= \frac{ (2m+1)^2 }{ 4n^2 }\\
		12n^2 &= (2m+1)^2
	\end{aligned}
\end{equation*}
On the left side we have an even integer and on the right side we have an odd integer---a contradiction. \contradiction

\item $a$ is an even integer and $b$ is an odd integer. Say $a = 2m$ and $b = 2n+1$, where $m$ and $n$ are integers. According to our assumption
\begin{equation*}
	\begin{aligned}
		\sqrt{3} &= \frac{ 2m }{ 2n+1 }\\
		3 &= \frac{ 4m^2 }{ (2n+1)^2 }\\
		3\ (2n+1)^2 &= 4m^2
	\end{aligned}
\end{equation*}
On the left side we have an odd integer and on the right side we have an even integer---a contradiction. \contradiction

\item $a$ and $b$ both are even integer. In this case, $a$ and $b$ must have 2 as a common factor. We can divide both $a$ and $b$ by 2. We can keep dividing by 2 as long as both remain even integer. At the end, we are in one of the previous three cases.
\end{enumerate}

\section*{Problem-229}
\subsection*{Problem Statement}
Prove that for $a \geq 0,\ b > 0$
\[
	\sqrt{ \frac{a}{b} } = \frac{ \sqrt{a} }{ \sqrt{b} }
\]

\subsection*{Solution}
Basically we are being asked to show that $\frac{ \sqrt{a} }{ \sqrt{b} }$ is the square root of the non-negative integer $\frac{a}{b}$. So, we need to prove the below two:
\begin{itemize}
	\item $\frac{ \sqrt{a} }{ \sqrt{b} }$ is non-negative.\\
	
	Since $a$ is non-negative, its square root must also be non-negative. Since $b$ is positive, its square root must also be positive. $\frac{ \sqrt{a} }{ \sqrt{b} }$ is thus the ratio of a non-negative number and a positive number; so, it must be non-negative.
	\item Squaring $\frac{ \sqrt{a} }{ \sqrt{b} }$ gives us $\frac{a}{b}$.

	\begin{equation*}
			\left( \frac{ \sqrt{a} }{ \sqrt{b} } \right)^2 = \frac{ (\sqrt{a})^2 }{ (\sqrt{b})^2 } = \frac{a}{b}.
	\end{equation*}		
\end{itemize}

\section*{Problem-230}
\subsection*{Problem Statement}
Is the equality $\sqrt{ a^2 } = a$ true for all $a$?

\subsection*{Solution}
It might seem counter-intuitive that the answer is no. But if we observe closely, the question is actually asking about the equality $+\sqrt{ a^2 } = a$, where the `+' is implied. When we make the `+' explicit, it is more noticeable that for $a < 0$, the equality cannot hold, because the left side of the equality is always positive. The  equality that works for all $a$ is $\sqrt{ a^2 } = |a|$, where
\[
	|a| = 
		\begin{cases}
		   \ \ a &\mbox{, if } a \geq 0\\
		   -a &\mbox{, if } a < 0
		\end{cases}
\]

It is unfortunate that we do not always make the `+' explicit.

\section*{Problem-232}
\subsection*{Problem Statement}
Which is bigger: $\sqrt{1001} - \sqrt{1000}$ or $\frac{1}{10}$?
\subsection*{Solution}
It would be easier to compare if we could convert $\sqrt{1001} - \sqrt{1000}$ into a fraction.
\begin{equation*}
	\begin{aligned}
		& \sqrt{1001} - \sqrt{1000}\\
		&= \frac{ \left(\sqrt{1001} - \sqrt{1000}\right)\ \left(\sqrt{1001} + \sqrt{1000}\right) }{ \sqrt{1001} + \sqrt{1000} }\\
		&= \frac{ \left( \sqrt{1001} \right)^2 - \left( \sqrt{1000} \right)^2 }{ \sqrt{1001} + \sqrt{1000} }\\
		&= \frac{1001 - 1000}{ \sqrt{1001} + \sqrt{1000} }\\
		&= \frac{1}{ \sqrt{1001} + \sqrt{1000} }\\
		&< \frac{1}{ \sqrt{900} + \sqrt{900} }
	\end{aligned}
\end{equation*}
Since $\sqrt{900} + \sqrt{900} = 60$, $\sqrt{1001} - \sqrt{1000} < \frac{1}{60} < \frac{1}{10}$. So, $\frac{1}{10}$ is the bigger between the two.

\section*{Problem-239}
\subsection*{Problem Statement}
\subsubsection{Solutions to quadratic equation}
A quadratic equation $x^2+px+q = 0$ can be solved by ``completing the square''.
\begin{equation*}
	\begin{aligned}
		x^2 + 2 \cdot x \cdot \frac{p}{2} + \left( \frac{p}{2} \right)^2 + \left( q - \frac{p^2}{4} \right) &= 0\\
		\left( x+\frac{p}{2} \right)^2 &= \frac{p^2}{4} - q\\
		x &= -\frac{p}{2} + \frac{ \pm \sqrt{p^2 - 4q} }{2}
	\end{aligned}
\end{equation*}
There are three possible cases:
\begin{enumerate}
	\item $p^2-4q > 0$. In this case two distinct solutions $x_1,\ x_2$ exist: 
	
	$x_1 = -\frac{p}{2} + \frac{ \sqrt{ p^2-4q } }{2}$ and $x_2 = -\frac{p}{2} - \frac{ \sqrt{ p^2-4q } }{2}$
	\item $p^2-4q = 0$. In this case the two solutions coincide; therefore, there is one distinct solution.
	
	$x_1 = -\frac{p}{2}$ and $x_2 = -\frac{p}{2}$
	\item $p^2-4q < 0$. In this case there is no real-valued solution.
\end{enumerate}
Since the sign of $p^2-4q$ determines how many distinct solutions there are, for convenience, we shall call this expression $D$, therefore, $D = p^2-4q$.

\subsubsection{Vieta's theorem}
If a quadratic equation $x^2 + px + q = 0$, has two distinct solutions $\alpha$ and $\beta$ then
\[
	\alpha + \beta = -p
\]
\[
	\alpha \cdot \beta = q
\]

\subsubsection{Proof-1 of Vieta's theorem}
If $x^2 + px + q = 0$ has two distinct solutions, we are in the first case of $0.0.1$; thus we can let $\alpha$ and $\beta$ be as follows:
\setcounter{equation}{0}
\begin{equation}
	\alpha = -\frac{p}{2} + \frac{ \sqrt{ D } }{2}
\end{equation}
\begin{equation}
	\beta = -\frac{p}{2} - \frac{ \sqrt{ D } }{2}
\end{equation}
$(1)+(2)$ gives:
\[
	\alpha + \beta = -p
\]
$(1) \times (2)$ gives:
\begin{equation*}
	\begin{aligned}
		\alpha \cdot \beta &= \frac{p^2}{4} - \frac{D}{4}\\
						   &= \frac{ p^2 - \left( p^2-4q \right) }{4}\\
						   &= q
	\end{aligned}
\end{equation*}

\subsubsection{Proof-2 of Vieta's theorem}
We can rewrite $x^2 + px + q = 0$ as $P(x) = 0$ where $P(x) = x^2 + px + q$ is a polynomial with degree two. If $\alpha$ and $\beta$ are two distinct solutions of $P(x) = 0$ then
\[
	P(x) = (x-\alpha)\ (x-\beta)\ R(x)
\]
Here $R(x)$ must be constant, otherwise right side would have degree more than two. In fact $R(x)$ must be exactly 1, otherwise the coefficient of $x^2$ would not match with the one in $P(x)$. So, we have
\begin{equation*}
	\begin{aligned}
		P(x) &= (x - \alpha)\ (x - \beta)\\
			 &= x^2 -(\alpha + \beta)\ x + \alpha \cdot \beta
	\end{aligned}
\end{equation*}
Comparing $P(x) = x^2 + (\ -(\alpha+\beta)\ )\ x + \alpha \cdot \beta$ with $P(x) = x^2 + px + q$ gives us Vieta's theorem.

\subsubsection{Questions to answer}
Can you generalize Vieta's theorem to the case of a quadratic equation having only one root? Are both proofs still valid for this case?

\subsection{Solution}
\subsubsection{Generalization}
Generalized version of Vieta's theorem is:\\

If a quadratic equation $x^2 + px + q = 0$ has two roots $\alpha$ and $\beta$ (the roots can be equal) then
\[
	\alpha + \beta = -p
\]
\[
	\alpha \cdot \beta = q
\]

\subsubsection{Proof-1}
Vieta's generalized theorem already works for the first case of $0.0.1$. We shall show that it works for the second case as well. In the second case of $0.0.1$, $\alpha = \beta = -\frac{p}{2}$.
\[
	\alpha + \beta = 2\alpha = 2 \left( -\frac{p}{2} \right) = -p
\]
In the second case of $0.0.1$ we also have $p^2-4q = 0$; thus $q = \frac{p^2}{4}$.
\[
	\alpha \cdot \beta = \alpha^2 = \frac{p^2}{4} = q
\]

\subsubsection{Proof-2}
Since $x^2+px+q = 0$ has two duplicate solutions, we can consider $P(x) = x^2 + px + q$ having two repeated roots, each equal to $\alpha$. Thus $P(x) = (x-\alpha)\ (x-\alpha)$. In other words, $P(x) = x^2 - 2\alpha + \alpha^2$. A comparison gives us Vieta's theorem.

\end{document}





















